\section{问题背景与重述}

\subsection{赛题背景}

电动汽车正在重塑交通的未来。电力公司需要预测电动汽车接入电网后的负荷，金融机构需要评估用户的用车经济压力，政策制定者则需要量化用户从燃油车转向电动车后所节省的潜在支出。这三类诉求分别对应消费者转型过程中的三个关键经济变量：月度能源消耗量、月度充电成本与燃油替代支出。《电动汽车革命：消费者行为与充电洞察》数据集收录 50,000 条消费者记录，涵盖人口统计、通勤行为、充电基础设施、环保意识、技术接受度及财务等维度，为上述三类预测提供了信息基础。

\subsection{三个子问题}

赛题的核心是三个回归任务。它们的预测目标与评分指标汇总在表\ref{tab:tasks}中。

\begin{table}[H]
\centering
\caption{赛题子问题概览}
\label{tab:tasks}
\begin{tabular}{cccc}
\toprule
\textbf{子问题} & \textbf{预测目标} & \textbf{类型} & \textbf{评分指标} \\
\midrule
任务一：月度能源消耗量预测 & monthly\_energy\_consumption\_kwh & 回归 & adj-$R^2$ \\
任务二：月度充电成本预测 & monthly\_charging\_cost & 回归 & adj-$R^2$ \\
任务三：燃油替代支出预测 & fuel\_expense\_per\_month & 回归 & adj-$R^2$ \\
\bottomrule
\end{tabular}
\end{table}

三个子问题共用同一份数据，但目标之间存在级联关系：月度充电成本在物理上近似等于月度能耗与电价的乘积，因此任务二的建模结果直接依赖任务一的预测精度。此外，任务三要求基于用户当前的出行习惯与车辆属性，预测其若继续驾驶燃油车的月度燃油支出，属于反事实推断。

\subsection{总体技术路线}

为保证结果可复现且不同方案之间具备可比性，三项任务共用一套实验框架：

\begin{enumerate}[itemsep=0pt]
\item 首先进行数据 Schema 校验与固定折划分：使用固定随机种子生成 5 折划分，并额外生成 3 组随机种子（42/2026/7）用于最终确认；
\item 所有数据变换（缺失值填补、编码器、缩尾边界）仅在训练折内部拟合，禁止利用验证折或测试集统计量；
\item 每项任务先以可解释的物理公式或线性模型建立稀疏基线，再通过特征块消融与残差诊断判断是否引入非线性模型；
\item 复杂模型仅在越过预设晋级门槛（OOF adj-$R^2$ 增益 $\ge 0.0002$ 且 Bootstrap 95\% 区间下界大于 0）时采用；
\item 模型确定后，在 5 折 $\times$ 3 种子上仅做一次最终确认，不再反复调整。
\end{enumerate}

本研究的核心原则是``物理结构优先、复杂模型克制''：当数据本身可由近似确定性的公式刻画时，稀疏模型不仅可解释性更强，且在调整决定系数的评分口径下更具优势——因为每增加一个无增益的预测变量都会带来自由度惩罚。
